% preamble.tex — en-tête LaTeX pour la sortie PDF (moteur LuaLaTeX)
% Quarto charge déjà booktabs, longtable, etc. ; on ajoute quelques réglages sûrs.

\usepackage{microtype}
\usepackage{amsmath}
\usepackage{amssymb}
\usepackage{eso-pic}

% Légendes compactes ; le caractère % dans une légende doit être écrit \%
\usepackage{caption}
\captionsetup{font=small, labelfont=bf}

% Interligne légèrement aéré
\usepackage{setspace}
\setstretch{1.05}

% Rappels (voir la préface) :
%  - éviter les indices Unicode (₁, ₒ) dans le texte ou les figures destinés au PDF ;
%    écrire les indices en LaTeX, p. ex. $Y_{\text{obs}}$, $t_1$.
%  - échapper le pourcentage dans les légendes et le texte : 30\%.

% Encadrés pédagogiques (filtre encadres.lua) : « À l'application » et « Ce qu'on doit voir »
\usepackage[most]{tcolorbox}
\definecolor{appfond}{HTML}{E6F6F4}\definecolor{appbord}{HTML}{17A398}
\definecolor{voirfond}{HTML}{FDF1E5}\definecolor{voirbord}{HTML}{D95F0E}
\newtcolorbox{encadreapp}{enhanced, breakable, frame hidden, boxrule=0pt, sharp corners,
  colback=appfond, borderline west={3pt}{0pt}{appbord}, left=8pt, right=6pt, top=4pt, bottom=4pt,
  before skip=8pt, after skip=8pt}
\newtcolorbox{encadrevoir}{enhanced, breakable, frame hidden, boxrule=0pt, sharp corners,
  colback=voirfond, borderline west={3pt}{0pt}{voirbord}, left=8pt, right=6pt, top=4pt, bottom=4pt,
  before skip=8pt, after skip=8pt}
